\section{参数校准与模型检验}

\subsection{模型校准与拟合表现}

\subsubsection{多目标校准}

本文采用多目标校准函数，同时约束总体 RMSE、峰值、末日价格和局部阶段解释能力。设 \(\theta\) 表示需要校准的参数向量，\(\Theta\) 为由赛题范围、物理上限和行为参数边界共同构成的可行域，则参数估计写为

\begin{equation}
  \hat\theta
  =
  \arg\min_{\theta\in\Theta}
  \mathcal L(\theta).
\end{equation}

其中 \(\Theta\) 对 SPR 上限、绕道能力、库存日释放、供应中断范围和弹性区间施加硬约束；行为参数只允许在经济含义一致的范围内搜索。目标函数为

\begin{equation}
\begin{aligned}
\mathcal L(\theta)={}&
\mathrm{RMSE}
+0.20|\Delta P_{\mathrm{peak}}|
+0.25|\Delta P_{\mathrm{final}}| \\
&+0.15\mathrm{RMSE}_{\mathrm{high}}
+0.12\mathrm{RMSE}_{\mathrm{early}}
+0.12\mathrm{RMSE}_{\mathrm{mid}} \\
&+0.18\mathrm{RMSE}_{\mathrm{late}}
+0.18\mathrm{RMSE}_{\mathrm{low}} .
\end{aligned}
\end{equation}

这里 \(\Delta P_{\mathrm{peak}}\) 与 \(\Delta P_{\mathrm{final}}\) 分别表示模拟峰值和期末价格相对附件真实价格的偏差；分段 RMSE 用于避免仅降低全窗口误差而牺牲前期冲击、高位平台或后期再定价解释能力。

校准过程包括 36000 组固定种子随机搜索、差分进化连续精修、800 组局部稳健性复核以及 60000 组拟合质量精修。所有精修均在既有模型结构内完成，不新增额外机制项。多组候选参数仅用于筛选稳定解，正文不展开逐一比较。

\begin{table}[htbp]
\centering
\caption{短期冲击模型综合最优参数}
\begin{threeparttable}
\begin{tabular}{l S[table-format=4.3] l S[table-format=3.3]}
\toprule
参数 & {数值} & 参数 & {数值} \\
\midrule
供应中断量 & 1493 & SPR 释放上限 & 670 \\
SPR 启动延迟 & 3 & 绕道启动日 & 8 \\
绕道能力上限 & 237 & 长期需求弹性 & -0.25 \\
恐慌初始强度 & 0.12 & 恐慌衰减速度 & 0.11 \\
库存日缓冲上限 & 250 & 价格传导系数 $\phi$ & 0.028 \\
调整速度 $v$ & 0.344 & 风险权重 $\beta$ & 2.462 \\
不确定性/制度风险强度 & 0.247 & 库存响应系数 & 0.666 \\
缓冲确认强度 & 0.220 & 预期修复强度 & 0.240 \\
\bottomrule
\end{tabular}
\begin{tablenotes}
\footnotesize
\item 注：供应中断量、SPR 释放上限、绕道能力上限和库存日缓冲上限的单位为万桶/日；其余无量纲参数用于刻画行为强度、调整速度或弹性。
\end{tablenotes}
\end{threeparttable}
\end{table}

表中参数可分为两类理解：供应中断量、SPR 上限、绕道能力、库存缓冲和需求弹性受题面范围约束；风险权重、价格调整速度、缓冲确认和预期修复强度属于行为价格形成参数，必须依赖数据校准和鲁棒性检验。本文后续所有模型评价均围绕这一区分展开，校准参数只在相应机制语境下解释。

其中，供应中断量 1493 万桶/日是在赛题给定中断范围内由短期冲突窗口校准得到的中心情景值，不等同于官方观测到的实际中断量。本文将其作为模型参数使用，并在长期乐观、中性、悲观情景以及蒙特卡洛抽样中覆盖 1400--1800 万桶/日的扰动区间。

\subsubsection{拟合结果}

短期冲击模型在冲突窗口内取得 RMSE 为 3.38 美元/桶、MAE 为 2.76 美元/桶。模拟峰值为 110.32 美元/桶，实际最高收盘价为 114.06 美元/桶，峰值低估 3.74 美元/桶；末日模拟价格为 110.32 美元/桶，实际末日价格为 110.31 美元/桶，末日误差仅为 0.01 美元/桶。在此基础上，本文只在高位平台形成和中期再定价回落两个薄弱窗口启用低自由度 Ridge 收益率修正，得到短期冲击模型增强路径，RMSE 降至 3.18 美元/桶、MAE 降至 2.69 美元/桶、MAPE 为 2.69\%，方向命中率为 68.9\%。

\keypoint{短期模型的核心价值在于解释极端冲突窗口中的前期冲击、中期平台和后期再定价；最终采用的阶段 Ridge 增强路径进一步修正平台期日度震荡。}

\begin{figure}[htbp]
\centering
\includegraphics[width=0.92\linewidth]{短期模型拟合效果.png}
\caption{短期冲击模型拟合效果}
\label{fig:short-term-fit}
\end{figure}

\begin{table}[htbp]
\centering
\caption{短期冲击模型分段误差}
\begin{tabular}{l S[table-format=2.0] S[table-format=1.2] S[table-format=1.2] S[table-format=1.2]}
\toprule
分段 & {样本数} & {RMSE} & {MAE} & {最大绝对误差} \\
\midrule
全窗口 & 46 & 3.38 & 2.76 & 7.24 \\
前期冲击 & 11 & 3.17 & 2.73 & 5.79 \\
中期平台 & 14 & 4.21 & 3.52 & 7.24 \\
后期再定价 & 21 & 2.82 & 2.27 & 5.37 \\
高价平台 & 4 & 2.70 & 2.19 & 3.93 \\
低价回落 & 21 & 3.20 & 2.76 & 5.79 \\
\bottomrule
\end{tabular}
\end{table}

\begin{figure}[htbp]
\centering
\includegraphics[width=0.92\linewidth]{短期模型机制贡献.png}
\caption{短期模型机制贡献分解}
\end{figure}

\subsection{基准对比与统计检验}

\subsubsection{基准对比与滞后检验}

金融时间序列中，上一日价格朴素基准是最重要参照。本文同时使用滚动 ARIMA(1,1,0) 作基础时间序列对照，所有基准在预测第 $t$ 日时只使用 $t-1$ 日及以前信息。三日均值和漂移随机游走也在程序中复核，结论与朴素基准一致，正文不再展开。

\begin{figure}[htbp]
\centering
\begin{tikzpicture}[
  node distance=0.7cm and 0.75cm,
  center/.style={draw=OilBorder, fill=PaperNoteBack, rounded corners=1pt, align=center, minimum height=0.78cm, text width=2.65cm, font=\small\bfseries},
  test/.style={draw=OilBorder, fill=white, rounded corners=1pt, align=center, minimum height=0.62cm, text width=2.35cm, font=\scriptsize},
  arrow/.style={-{Latex[length=2.0mm]}, line width=0.42pt, draw=OilGray}
]
\node[center] (model) {短期冲击模型};
\node[test, above left=of model] (baseline) {朴素基准\\ARIMA};
\node[test, above=of model] (lag) {滞后平移\\拐点检验};
\node[test, above right=of model] (ablation) {机制消融};
\node[test, below left=of model] (param) {参数扰动\\剖面检验};
\node[test, below=of model] (resid) {残差诊断\\DM 检验};
\node[test, below right=of model] (boundary) {历史窗口\\适用边界};
\foreach \x in {baseline,lag,ablation,param,resid,boundary}
  \draw[arrow] (\x) -- (model);
\end{tikzpicture}
\caption{模型鲁棒性检验框架}
\end{figure}

\begin{table}[htbp]
\centering
\caption{短期最终模型与关键基准对比}
\begin{threeparttable}
\begin{tabular}{l S[table-format=1.2] S[table-format=1.2] S[table-format=1.2] S[table-format=2.1]}
\toprule
模型 & {RMSE} & {MAE} & {MAPE(\%)} & {方向命中率(\%)} \\
\midrule
短期冲击模型 & 3.38 & 2.76 & 2.76 & 66.7 \\
\textbf{短期冲击模型增强路径} & \textbf{3.18} & \textbf{2.69} & \textbf{2.69} & \textbf{68.9} \\
朴素上一日基准 & 4.47 & 3.57 & 3.54 & 46.7 \\
滚动 ARIMA(1,1,0) 基准 & 4.53 & 3.67 & 3.63 & 46.7 \\
\bottomrule
\end{tabular}
\begin{tablenotes}
\footnotesize
\item 注：RMSE 与 MAE 单位为美元/桶；MAPE 与方向命中率以百分比计。短期冲击模型增强路径指“短期冲击模型 + 阶段 Ridge 修正”，只在高位平台形成和中期再定价回落两个残差薄弱段启用收益率修正，其他事件段保持机制路径不变。
\end{tablenotes}
\end{threeparttable}
\end{table}

\begin{figure}[htbp]
\centering
\includegraphics[width=0.86\linewidth]{短期模型基准对比.png}
\caption{短期冲击模型与关键基准 RMSE 对比}
\end{figure}

\cref{fig:short-term-fit} 展示了短期主模型的路径拟合结果。本文进一步将模型曲线整体向左、向右平移 1--5 天后重新计算 RMSE。结果显示，原始位置 RMSE 最低，为 3.38；向左平移 1 天后 RMSE 升至 3.90，向左平移 2 天后升至 4.71，说明模型拟合位置与真实冲击窗口基本同步。

\begin{figure}[htbp]
\centering
\includegraphics[width=0.84\linewidth]{短期模型滞后平移检验.png}
\caption{模型曲线平移后的 RMSE 检验}
\end{figure}

\subsection{模型鲁棒性与稳定性评估}

\subsubsection{结构稳定性检验}

由于冲突窗口只有 46 个交易日，而校准参数较多，本文进一步进行参数扰动与结构稳定性检验。结果显示，在 1000 组 $\pm 15\%$ 参数扰动中，96.6\% 的样本模拟峰值仍落在 105--125 美元/桶区间。本文因此将模型定位为“事件机制递推模型”，重点解释霍尔木兹封锁情境下的条件价格路径。

\begin{figure}[htbp]
\centering
\includegraphics[width=0.88\linewidth]{短期模型过拟合压力测试.png}
\caption{短期模型参数扰动检验}
\end{figure}

\subsubsection{分段误差与残差诊断}

仅报告全窗口 RMSE 容易掩盖局部失败。本文进一步检查逐日残差、滚动误差、实际值与模拟值一致性以及残差自相关。分段误差显示，中期平台是最难解释的部分，RMSE 为 4.39 美元/桶；后期再定价阶段 RMSE 降至 2.81 美元/桶，说明缓冲确认折价和预期修复折价对解释后期回落有实际作用。

\begin{figure}[htbp]
\centering
\includegraphics[width=0.92\linewidth]{短期模型残差诊断增强.png}
\caption{短期模型增强残差诊断}
\end{figure}

增强诊断表明，模型没有通过牺牲某一局部阶段来换取总体 RMSE。实际值与模拟值的一致性图说明模型能够跟随冲突窗口的主要价格平台；残差分布没有出现单侧系统性偏移；滚动误差在中期平台阶段略高，符合市场对封锁持续时间和缓冲机制可持续性反复定价的事实。残差一阶自相关约为 0.254，说明仍存在一定连续偏差，但未形成明显强自相关。

\begin{figure}[htbp]
\centering
\includegraphics[width=0.90\linewidth]{短期残差分段分布图.png}
\caption{短期模型残差的分段分布}
\label{fig:segmented-residual-distribution}
\end{figure}

进一步将逐日误差、事件段误差和机制项放在同一口径下复核，结果见\cref{fig:short-error-attribution}。最大单日误差出现在 2026 年 4 月 6 日，短期冲击模型低估 7.24 美元/桶；事件段层面，高位平台形成和中期再定价回落仍是主要难点，RMSE 分别约为 3.73 和 3.78 美元/桶。增强路径只针对这两个薄弱窗口加入低自由度 Ridge 收益率修正，使全窗口 RMSE 从 3.38 降至 3.18，方向命中率从 66.7\% 提高至 68.9\%。因此，本文报告的短期最优拟合效果来自“机制解释 + 局部收益惯性修正”，而不是多个候选模型的简单堆叠。

\begin{figure}[htbp]
\centering
\includegraphics[width=0.94\linewidth]{短期模型误差归因.png}
\caption{短期模型误差归因复核}
\label{fig:short-error-attribution}
\end{figure}

\subsubsection{时间轴语义与删块稳健性复核}

短期递推中的 SPR 释放、绕道恢复、需求收缩和恐慌衰减都依赖时间变量。若将每一行交易日都视为等间隔一步，可能低估周末和节假日期间政策与运输机制的实际推进；若采用日历日差，则相邻交易行在周末后会出现 3 天跳变。为避免时间轴选择成为隐藏建模假设，本文在不重新调参的条件下，对同一组最优参数分别采用日历日口径和交易日序号口径进行复算。

\begin{table}[htbp]
\centering
\caption{短期模型时间轴口径对比}
\begin{tabular}{lrrrr}
\toprule
时间轴口径 & 最大时间索引 & RMSE & MAE & 校准损失 \\
\midrule
日历日口径 & 64 & 3.380 & 2.756 & 6.505 \\
交易日序号口径 & 45 & 7.407 & 5.984 & 20.178 \\
\bottomrule
\end{tabular}
\end{table}

结果显示，交易日序号口径在同参条件下明显恶化，尤其高价平台 RMSE 升至 15.22 美元/桶。因此，当前主模型保留日历日口径是合理的：封锁冲击、政策释放和运输恢复并不会在周末停止，按真实经过时间推进更符合机制解释题的语义。

进一步地，本文将 46 个交易日按时间顺序切分为 5 个连续块，使用同一组参数检查各块误差。被删除块 RMSE 分别为 3.27、3.58、4.60、2.55 和 2.84 美元/桶，标准差为 0.71。最难片段为 2026 年 3 月 27 日至 4 月 9 日，RMSE 为 4.60 美元/桶，说明模型的主要残余误差集中在中期平台再定价阶段，而不是平均分布在所有时期。

\begin{figure}[htbp]
\centering
\includegraphics[width=0.92\linewidth]{短期模型时间轴与删块稳健性.png}
\caption{短期模型时间轴与删块稳健性复核}
\end{figure}

\subsubsection{统计检验与 R 语言辅助定位}

本文还在剔除首日后构造可比较样本，将短期冲击模型与朴素上一日基准进行误差和统计比较。剔除首日是因为朴素上一日基准需要上一交易日真实价格作为输入，因此本节统计检验中的 RMSE 与前文全窗口 RMSE 略有差异。

\begin{table}[htbp]
\centering
\caption{短期模型统计检验结果}
\begin{tabularx}{\linewidth}{lrrX}
\toprule
检验项 & 本文模型 & 对照或 p 值 & 解释 \\
\midrule
RMSE & 3.42 & 4.52 & 对照为朴素上一日基准，改善约 24.4\% \\
MAE & 2.80 & 3.65 & 平均绝对误差低于朴素基准 \\
MAPE & 2.79\% & 3.62\% & 相对误差低于朴素基准 \\
DM 平方损失检验 & -2.12 & 0.020 & R 独立复核下，单侧检验支持模型优于朴素基准 \\
DM 绝对损失检验 & -1.82 & 0.038 & 绝对损失优势接近 5\% 显著性边界，不作强显著性夸大 \\
方向命中率 & 65.9\% & 0.024 & 相对随机 50\% 的单侧二项检验 \\
Ljung--Box lag10 & 10.92 & 0.364 & 未发现显著线性残差自相关 \\
ARCH--LM lag5 & 5.22 & 0.389 & 未发现强条件异方差证据 \\
Newey--West 校准诊断 & 1.061 & $<0.001$ & 仅作拟合值与真实价格的刻画性校准诊断，不作为独立因果或样本外证明 \\
\bottomrule
\end{tabularx}
\end{table}

\begin{figure}[htbp]
\centering
\includegraphics[width=0.92\linewidth]{短期模型统计审计.png}
\caption{短期模型统计检验图}
\end{figure}

分段统计给出更细的判断：前段冲击吸收和后段预期修复阶段，本文模型相对朴素基准的 RMSE 改善率分别约为 33.4\% 和 34.1\%；中段平台维持阶段改善率约为 3.9\%，说明平台震荡期仍是模型的主要难点。总体看，DM 检验和方向命中检验提供了支持性统计证据；样本只有 45 个可比较交易日，绝对损失 DM 检验也接近显著性边界，因此这些统计结果作为基准对比、滞后检验和机制消融之外的补充验证。Newey--West 回归仅用于刻画拟合路径与真实价格的线性一致性。

R 语言在本文中承担独立计量复核功能，用于检验 DM、方向命中、Ljung--Box、ADF、ARCH/GARCH 和 Newey--West 稳健校准回归。R 复核显示，价格水平和对数价格不能拒绝单位根假设，而对数收益率可拒绝单位根假设；GARCH(1,1) 在短样本下提示波动持续性。除计量检验外，本文还使用 R 的 ggplot2 对关键证据图进行学术化重绘。

\begin{figure}[htbp]
\centering
\includegraphics[width=0.92\linewidth]{R短期误差学术诊断.png}
\caption{R/ggplot2 短期误差学术诊断}
\end{figure}

\subsection{机制消融、事件边界与机器学习辅助}

\subsubsection{机制消融的定量证据}

本文对完整模型进行逐项消融：关闭某一机制时，其他参数保持不变，不重新校准。该设计用于检验单个机制在同一参数体系下的边际贡献，而不是重新给每个缺失机制模型寻找最优参数。

\begin{table}[htbp]
\centering
\caption{短期模型机制消融实验}
\begin{tabular}{lrrrr}
\toprule
消融实验 & RMSE & RMSE变化 & 后期RMSE & 低价回落RMSE \\
\midrule
完整模型 & 3.38 & 0.00 & 2.82 & 3.20 \\
无 SPR & 5.19 & +1.81 & 5.08 & 5.90 \\
无库存 & 4.08 & +0.70 & 2.63 & 4.39 \\
无绕道 & 3.38 & -0.00 & 2.83 & 3.22 \\
无需求收缩 & 4.51 & +1.13 & 3.54 & 5.17 \\
无恐慌溢价 & 3.67 & +0.29 & 2.80 & 3.39 \\
无风险溢价 & 30.34 & +26.96 & 35.50 & 25.57 \\
无缓冲折价 & 4.55 & +1.17 & 3.77 & 5.38 \\
无修复折价 & 7.13 & +3.75 & 9.83 & 9.51 \\
无市场修正 & 27.02 & +23.64 & 29.32 & 18.92 \\
\bottomrule
\end{tabular}
\end{table}

消融结果可以分成两层理解。第一层是题面物理层：SPR 释放、商业库存和需求收缩的影响最清楚，去除 SPR 后 RMSE 从 3.38 升至 5.19，去除商业库存后 RMSE 升至 4.08，去除需求收缩后低价回落 RMSE 升至 5.17，说明政策缓冲、库存吸收和需求侧调整不是装饰项。绕道运输在 46 个交易日短窗口内边际影响较小，但它在中长期油价调节模型中会影响剩余缺口和二次跳涨风险，因此不应因为短窗口贡献小就删除。

第二层是市场价格形成层：风险与不确定性溢价负责解释价格快速进入高位平台，缓冲确认折价和预期修复折价负责解释后期再定价。若去除预期修复折价，后期 RMSE 从 2.82 升至 9.83，说明市场对政策缓冲和运输恢复的预期并不是可有可无的修饰，而是解释后期回落的关键机制。

\subsubsection{阶段 Ridge 辅助边界}

历史价格特征并不能单独解释霍尔木兹封锁冲突窗口。本文使用冲突前历史滞后价格特征训练低自由度 Ridge 收益率模型，目的只是刻画价格自身的短期惯性，并将其限制在两个残差薄弱窗口内使用。程序中还比较了纯 Ridge、全窗口 Ridge 和在线残差校正等方案；这些方案或解释性不足，或会损害初始冲击与后期二次抬升阶段，因此不作为正文主线展开。最终采用的阶段 Ridge 辅助层只修正高位平台和中期再定价中的日度收益惯性，主解释仍由供应缺口、SPR、库存、绕道运输、恐慌溢价和预期修复机制承担。

\subsubsection{拐点、预测步长与事件边界}

本文模型是条件机制递推模型。模型给定冲突开始、题面参数和校准参数后，从初始价格开始逐日递推；真实价格用于参数校准和事后评价，不在每日递推时作为当日输入。因此，RMSE 为 3.38 的含义是模型对整段冲突路径的条件解释误差。

为了检查模型在突变日是否只是慢半拍跟随真实价格，本文选取实际价格绝对变动最大的 10 个交易日进行方向检查。结果显示，模型在 8 个交易日与真实价格同向同步变化，2 个交易日体现为提前 1 日变化。这说明模型对主要冲击方向具有同步解释力，但个别日度噪声仍不能完全捕捉。

\begin{figure}[htbp]
\centering
\includegraphics[width=0.92\linewidth]{短期模型拐点局部检验.png}
\caption{短期模型拐点局部检验}
\end{figure}

进一步地，本文将短期冲突窗口拆成四个事件段，分别检验初始冲击吸收、高位平台形成、中期再定价回落和后期二次抬升。\cref{tab:short-event-window-audit} 显示，四个事件窗口下模型 RMSE 均低于朴素上一日基准，其中初始冲击吸收和后期二次抬升窗口的相对改善率超过 30\%。高位平台形成和中期再定价回落窗口误差相对更高，说明模型最难处理的不是方向判断，而是平台期真实价格的短期反复震荡。

\begin{table}[htbp]
\centering
\caption{短期模型事件窗口检验表}
\label{tab:short-event-window-audit}
\begin{tabular}{lrrrr}
\toprule
事件窗口 & RMSE & MAE & 方向命中率 & 相对朴素改善率 \\
\midrule
初始冲击吸收 & 3.27 & 2.88 & 77.8\% & 32.7\% \\
高位平台形成 & 3.91 & 3.23 & 70.0\% & 16.5\% \\
中期再定价回落 & 3.90 & 3.35 & 63.6\% & 14.1\% \\
后期二次抬升 & 2.66 & 2.14 & 54.5\% & 32.0\% \\
\bottomrule
\end{tabular}
\end{table}

\begin{figure}[htbp]
\centering
\includegraphics[width=0.94\linewidth]{短期模型事件窗口检验.png}
\caption{短期模型事件窗口检验}
\end{figure}

本文还进行了历史高波动窗口边界检验：固定霍尔木兹机制参数，直接套用到 2020 年、2022 年等非霍尔木兹事件窗口。结果显示，该模型在这些窗口通常显著弱于朴素上一日基准。这说明本模型高度特化于霍尔木兹封锁事件下的系统机制解释，其结论应限定在极端地缘冲突下的条件路径解释与情景预测。

\subsubsection{参数自由度与稳健性区间}

短期冲突窗口只有 46 个交易日，而连续校准参数共有 21 个，其中 10 个属于题面或机制范围内的物理/半物理参数，另外 11 个属于行为与价格调整参数。为刻画参数自由度带来的不确定性，本文进一步报告压力测试、参数剖面和局部扰动结果。

\begin{table}[htbp]
\centering
\caption{\(\pm 15\%\) 参数压力测试摘要}
\begin{tabular}{lrrrr}
\toprule
指标 & 均值 & 中位数 & 90\%分位数 & 最大值 \\
\midrule
RMSE & 4.96 & 4.74 & 6.33 & 10.27 \\
模拟峰值 & 110.39 & 110.30 & 114.16 & 119.62 \\
模拟末日价格 & 110.39 & 110.30 & 114.16 & 119.62 \\
高价平台 RMSE & 3.90 & 3.36 & 6.84 & 11.84 \\
低价回落 RMSE & 4.98 & 4.63 & 7.71 & 11.60 \\
\bottomrule
\end{tabular}
\end{table}

压力测试显示，在 1000 组 \(\pm 15\%\) 参数扰动中，61.0\% 的样本 RMSE 不高于 5，96.6\% 的样本模拟峰值落在 105--125 美元/桶区间，全部样本末日价格落在 100--120 美元/桶区间。这个结果意味着，价格平台并不依赖唯一一组精确参数点；但较大扰动会削弱 RMSE 表现，因此本文只主张模型具有结构性稳健性，不主张任意参数下都能保持高精度。

本文固定其他参数，只对价格传导系数、封锁风险权重、不确定性底座、价格调整速度、缓冲确认折价和预期修复折价进行一维剖面检验。\cref{fig:short-parameter-profile} 表明，当前最优参数附近的误差曲线整体连续；其中封锁风险权重和不确定性底座对 RMSE 更敏感，说明风险溢价通道是短期路径解释的关键部分。

\begin{figure}[htbp]
\centering
\includegraphics[width=0.94\linewidth]{短期模型参数剖面图.png}
\caption{短期模型关键参数一维剖面}
\label{fig:short-parameter-profile}
\end{figure}

\begin{figure}[htbp]
\centering
\includegraphics[width=0.92\linewidth]{短期模型稳健性带.png}
\caption{短期模型局部扰动稳健性区间}
\end{figure}

围绕当前入选参数进行的 800 组局部扰动进一步显示，在 801 组样本中，99.4\% 的样本同时满足 RMSE 不高于 5、高价平台 RMSE 不高于 5、低价回落 RMSE 不高于 6、后期 RMSE 不高于 5 的检验阈值。通过阈值样本的 10\%--90\% 价格带平均宽度为 2.18 美元/桶，最大宽度为 3.01 美元/桶，说明当前参数附近存在稳定的可接受参数邻域。

针对多目标校准权重可能成为隐藏超参数的问题，本文进一步对已保存的前 10 个候选参数进行权重敏感性重排。将峰值误差、末日误差、分段 RMSE 等权重整体上调或下调 20\%，并分别强化峰值末日、分段误差和平台误差后，当前综合最优候选在 6/6 个权重情景下仍保持第一。该检验不等同于重新全局校准，但说明当前候选并非只依赖某一组精确权重偶然胜出。

本文将关键检验证据按“问题--证据--结论”的形式汇总于\cref{tab:short-validity-summary}。这些证据用于说明短期冲击模型是否能够支撑赛题结论，最终有效性由误差、机制解释和稳健性检验共同决定。

\begin{table}[htbp]
\centering
\caption{短期冲击模型有效性证据汇总}
\label{tab:short-validity-summary}
\begin{tabular}{p{0.23\linewidth}p{0.32\linewidth}p{0.33\linewidth}}
\toprule
检验问题 & 主要证据 & 论文中的作用 \\
\midrule
是否优于简单基准 & RMSE、MAE 和 MAPE 均优于朴素上一日、三日均值、漂移随机游走和滚动 ARIMA & 说明模型不是仅复制油价惯性 \\
是否存在滞后复制 & 最优平移仍为 0 日，拐点窗口中峰值与回落方向基本一致 & 回应“马后炮曲线”风险 \\
机制项是否必要 & 机制消融显示关闭风险、缓冲和预期修复通道会明显改变误差结构 & 说明扩展项具有边际解释贡献 \\
参数是否过度依赖单点 & 一维剖面连续，局部扰动样本形成窄价格带 & 支撑参数邻域稳定性 \\
残差是否仍有边界 & 残差存在一定波动聚集，方向命中率仍有提升空间 & 明确模型是事件机制解释模型，不是通用交易预测器 \\
\bottomrule
\end{tabular}
\end{table}

\subsubsection{内生性、反事实基准与结构边界}

模型中的恐慌初始强度和恐慌衰减速度用于刻画冲突初期风险溢价衰减，并非由新闻文本统计得到的外生恐慌指数。若后续加入新闻情绪模块，应使用 \(t-1\) 期或闭市前新闻解释 \(t\) 日价格，并进行格兰杰因果或滞后相关检验。

关于 200 美元以上反事实基准，传统供需模型的高价结果来自低短期弹性条件下的线性需求反事实基准，而不是现实预测；常弹性需求口径只作为机械上界参照。弹性敏感性检验显示，当需求价格弹性由 \(-0.05\) 放宽至 \(-0.25\) 或 \(-0.35\) 时，线性需求反事实价格会明显下降，部分情景已经接近或低于 120 美元/桶。因此，本文要论证的是：低短期弹性静态供需模型会显著高估现实价格，需要引入库存、SPR、绕道、需求收缩、价格传导折价和预期修复等动态机制解释现实平台。

递推结构检查显示，短期冲击模型函数没有设置 \(P_t>120\) 后强制压低价格的条件分支，也没有将 120 美元设置为价格上限。模型中的数值安全裁剪上限为 180 美元/桶，校准路径中从未触发；110--120 美元只作为题面观测参考区间出现在图表阴影和展示配置中，不参与价格递推计算。当前模型可解释 120 美元附近平台的机制来源，但不将其表述为多智能体动态纳什均衡。

\subsubsection{鲁棒性与边界检验总览}

本文将关键不确定性来源转化为鲁棒性与边界检验。\cref{tab:adversarial-tests} 中的检验均不重新寻找最优参数，而是在同一模型框架下改变某个关键条件，观察主结论是否保持稳定。

\begin{table}[htbp]
\centering
\caption{鲁棒性与边界检验总览}
\label{tab:adversarial-tests}
\begin{tabularx}{\linewidth}{lX X X}
\toprule
检验假设 & 检验方式 & 关键结果 & 论文含义 \\
\midrule
相对朴素基准的有效性 & 与上一日价格、三日均值、漂移随机游走和滚动 ARIMA 比较 & 可比样本 RMSE 为 3.42，朴素基准为 4.52；DM 平方损失单侧 \(p=0.020\) & 机制递推模型相对无变化预测具有误差优势 \\
滞后复制风险 & 平移检验、拐点局部检验和残差自相关检验 & Ljung--Box lag10 \(p=0.364\)，未发现显著线性残差自相关 & 拟合曲线没有明显依赖“昨天当今天”的滞后复制 \\
题面弹性与价格传导的区分 & 将价格传导系数强制设为 \(\phi=1\) & RMSE 升至 26.55，模拟峰值升至 166.84 美元/桶 & 题面低弹性用于反事实上界，\(\phi\) 是期货市场折价传导参数 \\
参数扰动稳定性 & 1000 组 \(\pm15\%\) 扰动和 800 组局部扰动 & 96.6\% 扰动样本峰值仍在 105--125 美元/桶，局部通过率 99.4\% & 价格平台不依赖唯一一组参数点 \\
长期点预测边界 & 蒙特卡洛 2000 条条件路径与三情景中心线对比 & 第 180 天 P05--P95 为 89.58--106.44 美元/桶，突破 120 概率约 16.3\% & 长期结论应读作条件区间和尾部风险，而非精确日度预测 \\
长期制度风险边界 & 将不确定性溢价减半或随时间强制衰减 & 中性第 180 天降至 89.15 或 87.82 美元/桶，悲观仍保留 99.80--101.27 美元/桶压力 & 制度风险影响长期平衡点，但悲观尾部仍来自多因素组合 \\
\bottomrule
\end{tabularx}
\end{table}
